\section{Il Codice CFD in Fortran~90}
\label{sec:codice_cfd}

\subsection{Equazioni di Governo: Eulero 2D}

Il codice risolve le \textbf{equazioni di Eulero 2D}, che descrivono il moto di un fluido
comprimibile non viscoso. Nella forma conservativa differenziale:

\begin{equation}
  \parziale{\uvec}{t}
  + \parziale{\fvec}{x}
  + \parziale{\gvec}{y} = \bm{0}
  \label{eq:eulero_diff}
\end{equation}

dove il vettore delle variabili conservative e i vettori dei flussi sono:

\begin{equation}
  \uvec = \begin{pmatrix} \rho E \\ \rho \\ \rho u \\ \rho v \end{pmatrix}, \quad
  \fvec = \begin{pmatrix} u(\rho E + P) \\ \rho u \\ \rho u^2 + P \\ \rho u v \end{pmatrix}, \quad
  \gvec = \begin{pmatrix} v(\rho E + P) \\ \rho v \\ \rho u v \\ \rho v^2 + P \end{pmatrix}
\end{equation}

con $\rho$ densità, $u, v$ componenti di velocità, $P$ pressione, $E$ energia totale specifica.
Il sistema è chiuso dall'equazione di stato dei gas perfetti:

\begin{equation}
  P = (\gamma - 1)\left(\rho E - \tfrac{1}{2}\rho(u^2+v^2)\right), \quad \gamma = 1.4
\end{equation}

\subsection{Adimensionalizzazione}

Le equazioni vengono risolte in forma adimensionale. Scelti $P_{\rm ref} = P^\circ_{\rm in}$
e $T_{\rm ref} = T^\circ_{\rm in}$ come grandezze di riferimento, le scale caratteristiche
derivate sono:
\begin{align*}
  \rho_{\rm ref} &= \frac{P_{\rm ref}}{R T_{\rm ref}}, &
  u_{\rm ref} &= \sqrt{R T_{\rm ref}}, &
  t_{\rm ref} &= \frac{L_{\rm ref}}{u_{\rm ref}}
\end{align*}

Una proprietà fondamentale dell'adimensionalizzazione è che le equazioni di Eulero mantengono
la stessa forma: i fattori di scala si elidono, e il codice lavora quindi con le stesse
espressioni indipendentemente dall'adimensionalizzazione scelta. In pratica, fissando
$P^\circ_{\rm in} = 1$ e $T^\circ_{\rm in} = 1$, le condizioni al contorno e le variabili
di campo assumono valori $O(1)$, migliorando il condizionamento numerico del sistema.

Le variabili adimensionalizzate rilevanti sono:
\begin{align}
  \bar{P} &= \frac{P}{P_{\rm ref}}, & \bar{T} &= \frac{T}{T_{\rm ref}}, &
  \bar{\rho} &= \frac{\bar{P}}{\bar{T}}, \\
  \bar{u} &= \frac{u}{u_{\rm ref}} = M \bar{a}, &
  \bar{a} &= \sqrt{\gamma \bar{T}}, &
  \bar{S} &= \gamma \ln \bar{T} - (\gamma-1)\ln \bar{P}
\end{align}

L'entropia adimensionale $\bar{S}$ è particolarmente utile per la validazione: in un campo
di Eulero privo di onde d'urto, $\bar{S}$ è conservata lungo le linee di flusso, quindi
ogni deviazione da zero è imputabile all'errore numerico di discretizzazione.

\subsection{Discretizzazione a Volumi Finiti}

Integrando l'equazione~\eqref{eq:eulero_diff} sul volume della cella $\Omega_i$:
\begin{equation}
  \parziale{}{t} \int_{\Omega_i} \uvec \, d\Omega
  + \oint_{\partial\Omega_i} \left(\fvec n_x + \gvec n_y\right) dS = \bm{0}
\end{equation}

Approssimando con il valore medio sulla cella e la somma dei flussi sui lati:
\begin{equation}
  \parziale{\uvec_i}{t} = -\frac{1}{|\Omega_i|}\sum_{j \in \mathcal{N}(i)}
    \Fvec_{ij} \cdot \nhat_{ij} \, l_{ij}
  \label{eq:volumi_finiti}
\end{equation}

dove $\Fvec_{ij}$ è il flusso numerico all'interfaccia tra la cella $i$ e il suo vicino $j$,
$\nhat_{ij}$ la normale uscente da $i$, e $l_{ij}$ la lunghezza dell'interfaccia.

\subsection{Calcolo dei Flussi: Schema di Lax--Friedrichs e Schema di Roe}

\subsubsection{Schema di Lax--Friedrichs Locale}

Lo schema di Lax--Friedrichs è uno schema \emph{centrato} con dissipazione numerica
proporzionale al massimo valore proprio locale:
\begin{equation}
  \Fvec_{ij} = \frac{1}{2}\left[F^L + F^R\right]
               - \frac{\lambda_{\max}}{2}\left(\uvec^R - \uvec^L\right)
\end{equation}
dove $\lambda_{\max} = \max(|\tilde{u}^L|+a^L,\,|\tilde{u}^R|+a^R)$ è la velocità massima
di propagazione e $\tilde{u} = \bm{q}\cdot\nhat$ è la velocità normale all'interfaccia.

I flussi vengono calcolati nel sistema di riferimento \emph{ruotato} locale all'interfaccia,
sfruttando la matrice di rotazione:
\begin{equation}
  \begin{pmatrix} \tilde{u} \\ \tilde{v} \end{pmatrix} =
  \begin{pmatrix} n_x & n_y \\ -n_y & n_x \end{pmatrix}
  \begin{pmatrix} u \\ v \end{pmatrix}
\end{equation}
e i flussi di quantità di moto vengono poi riportati nel sistema globale tramite la trasposta
della stessa matrice (che è ortogonale: $R^{-1} = R^T$).

\subsubsection{Schema di Roe (Upwind)}

Lo schema di Roe è un metodo \emph{upwind} che tiene conto della propagazione direzionale
delle onde caratteristiche. Il flusso è:
\begin{equation}
  \Fvec_{ij} = \frac{1}{2}\left[F^L + F^R\right]
               - \frac{1}{2} \sum_{k=1}^{4} |\hat{\lambda}_k| \hat{\alpha}_k \hat{r}_k
\end{equation}
dove $\hat{\lambda}_k$, $\hat{r}_k$, $\hat{\alpha}_k$ sono rispettivamente i valori propri,
i vettori propri destri e le \emph{wave strengths} della matrice Jacobiana di Roe, costruita
utilizzando le medie di Roe tra gli stati $L$ e $R$:
\begin{align}
  \sqrt{\hat{\rho}} &= \sqrt{\rho^L\rho^R}, &
  \hat{u} &= \frac{\sqrt{\rho^L}\,u^L + \sqrt{\rho^R}\,u^R}{\sqrt{\rho^L}+\sqrt{\rho^R}}, &
  \hat{H} &= \frac{\sqrt{\rho^L}\,H^L + \sqrt{\rho^R}\,H^R}{\sqrt{\rho^L}+\sqrt{\rho^R}}
\end{align}

Lo schema di Roe è più accurato di Lax--Friedrichs perché è meno dissipativo e cattura
meglio le discontinuità (onde d'urto) e le strutture fini del campo. Include la
\textit{entropy fix} di Harten per correggere le singolarità in corrispondenza dei punti sonici.

\subsection{Integrazione Temporale}

L'integrazione nel tempo è esplicita con schema di Eulero in avanti:
\begin{equation}
  \uvec_i^{k+1} = \uvec_i^k + \dt \parziale{\uvec_i}{t}^k
\end{equation}

Il passo temporale $\dt$ è calcolato ad ogni iterazione imponendo la condizione CFL:
\begin{equation}
  \dt = \text{CFL} \cdot \min_i \frac{\sqrt{|\Omega_i|}}{\lambda_{\max,i}}
  \label{eq:cfl}
\end{equation}

dove $\lambda_{\max,i} = |\bm{q}_i| + a_i$ è la velocità di propagazione massima nella
cella $i$. Il fattore $\sqrt{|\Omega_i|}$ approssima la dimensione caratteristica della cella
in 2D ($h \propto \sqrt{A}$, analogamente a $h \propto \sqrt[3]{V}$ in 3D).

\subsection{Condizioni al Contorno}

Le condizioni al contorno dipendono dalla direzione del flusso e dal regime:
\begin{itemize}
  \item \textbf{Inlet subsonico}: si impongono $T^\circ_{\rm in}$, $P^\circ_{\rm in}$,
    $\alpha_{\rm in}$; il numero di Mach è estrapolato dall'interno.
  \item \textbf{Inlet supersonico}: si impongono tutte e quattro le variabili conservative
    (flusso completamente entrante, nessuna caratteristica uscente).
  \item \textbf{Outlet subsonico}: si impone la pressione statica $P_{\rm exit}$;
    le altre tre variabili sono estrapolate dall'interno.
  \item \textbf{Outlet supersonico}: tutte le variabili sono estrapolate (nessuna
    caratteristica entrante).
  \item \textbf{Parete solida}: condizione di impermeabilità ($\bm{q}\cdot\nhat = 0$),
    implementata con una cella specchio che annulla la componente normale della velocità.
\end{itemize}

\subsection{Analisi di Convergenza della Griglia}

\subsubsection{Norma $L_2$ dell'Entropia}

Per quantificare l'errore numerico si usa la norma $L_2$ dell'entropia adimensionale $\bar{S}$,
che in un campo di Eulero senza onde d'urto dovrebbe essere identicamente nulla:
\begin{equation}
  \norma{\bar{S}}_2 = \sqrt{\frac{\displaystyle\sum_i \bar{S}_i^2 \cdot |\Omega_i|}
                                  {\displaystyle\sum_i |\Omega_i|}}
  \label{eq:norma_entropia}
\end{equation}

Questa norma è una misura dell'errore di discretizzazione: al raffinamento della griglia
($\lc \to 0$) ci si aspetta che $\norma{\bar{S}}_2 \to 0$.

\subsubsection{Ordine di Convergenza e Estrapolazione di Richardson}

Assumendo che l'errore di discretizzazione si comporti come:
\begin{equation}
  E = u_h - u_{\rm esatto} = k \cdot h^p + \mathcal{O}(h^{p+1})
\end{equation}
dove $h$ è la lunghezza caratteristica e $p$ l'ordine di convergenza, da due simulazioni con
$h_1$ e $h_2 = r h_1$ si ottiene:
\begin{equation}
  p = \frac{\ln(E_1/E_2)}{\ln(r)}, \quad
  u_{\rm esatto} \approx \frac{r^p u_{h_1} - u_{h_2}}{r^p - 1}
\end{equation}

Se la soluzione esatta non è nota a priori, si usa l'estrapolazione di Richardson con
\emph{ordine effettivo}: da tre simulazioni con $h$, $2h$, $4h$:
\begin{equation}
  p = \frac{\ln\!\left(\dfrac{u_{4h}-u_{2h}}{u_{2h}-u_h}\right)}{\ln 2}, \qquad
  u_{\rm esatto} \approx u_h + \frac{u_{2h}-u_h}{2^p - 1}
\end{equation}

Il \textbf{Grid Convergence Index} (GCI) fornisce una stima conservativa dell'incertezza
numerica:
\begin{equation}
  \text{GCI} = F_s \cdot |E_h|, \quad F_s = 3 \text{ (fattore di sicurezza)}
\end{equation}